% id: 39-22
% book: 베이직쎈 확률과 통계 · 03 확률의 뜻과 활용
% section: 개념 11 합사건, 곱사건, 배반사건, 여사건
% figure: none
% answer: 
% answer_source: 
% variant_level: 0 (원본 전사)
% 이 파일은 build.mjs 가 items.json 에서 만든다. 고칠 때는 items.json 을 고친다.
\smash{\raisebox{1.5mm}{\dmkichul{베이직쎈 확률과 통계 39쪽 22번}}}
표본공간 $S=\{1{,}\mkern3mu\mathopen{}2{,}\mkern3mu\mathopen{}3{,}\mkern3mu\mathopen{}\ldots{,}\mkern3mu\mathopen{}9\}$의 세 사건\par\smallskip\dispeq{$A=\{1{,}\mkern3mu\mathopen{}4{,}\mkern3mu\mathopen{}7{,}\mkern3mu\mathopen{}9\}$, $B=\{5{,}\mkern3mu\mathopen{}6{,}\mkern3mu\mathopen{}7{,}\mkern3mu\mathopen{}8\}$,}\par\dispeq{$C=\{2{,}\mkern3mu\mathopen{}3{,}\mkern3mu\mathopen{}6\}$}\par\smallskip\noindent 에 대하여 $A$와 $B$, $B$와 $C$, $C$와 $A$ 중 두 사건이 서로 배반사건인 것을 모두 찾으시오.
